\DocumentMetadata{
  pdfversion=2.0,
  pdfstandard=ua-2,
  lang=en-US,
  tagging=on,
  tagging-setup={math/setup=mathml-SE,tabsorder=structure}
}
\documentclass[11pt,letterpaper]{article}
\usepackage[margin=0.68in,includefoot]{geometry}
\usepackage{newcomputermodern}
\usepackage{amsmath,amscd,mathtools}
\usepackage{tikz,adjustbox}
\providecommand{\square}{\mdlgwhtsquare}
\usepackage[hidelinks]{hyperref}
\hypersetup{
  pdftitle={Topology PhD Exam, May 2017},
  pdfauthor={University of Florida Department of Mathematics},
  pdfsubject={Graduate mathematics examination},
  pdfdisplaydoctitle=true,
  bookmarksopen=true
}
\setcounter{secnumdepth}{0}
\setlength{\parindent}{0pt}
\setlength{\parskip}{0.28em}
\setlength{\emergencystretch}{3em}
\setlength{\leftmargini}{2.15em}
\setlength{\leftmarginii}{2.65em}
\setlength{\leftmarginiii}{3em}
\renewcommand{\labelenumi}{\textbf{\theenumi.}}
\pagestyle{empty}
\raggedbottom
\allowdisplaybreaks
\newenvironment{examproblems}[1]{%
  \begin{enumerate}%
  \setcounter{enumi}{#1}%
  \setlength{\topsep}{0.35em}%
  \setlength{\partopsep}{0pt}%
  \setlength{\itemsep}{0.48em}%
  \setlength{\parsep}{0.18em}%
}{\end{enumerate}}
\newenvironment{examsubparts}{%
  \begin{enumerate}%
  \setlength{\topsep}{0.18em}%
  \setlength{\partopsep}{0pt}%
  \setlength{\itemsep}{0.16em}%
  \setlength{\parsep}{0.1em}%
}{\end{enumerate}}
\newenvironment{examinstructions}{%
  \par\begingroup\itshape
}{%
  \par\endgroup\vspace{0.18em}
}
\makeatletter
\renewcommand\section{\@startsection{section}{1}{0pt}%
  {-1.25ex plus -.25ex minus -.1ex}%
  {0.55ex plus .1ex}%
  {\normalfont\large\bfseries}}
\renewcommand{\@maketitle}{%
  \newpage\null\vskip -1em%
  {\footnotesize\bfseries
    \href{https://www.ufl.edu/}{University of Florida}, \href{https://math.ufl.edu/}{Department of Mathematics}\par}%
  \vskip -0.5em%
  \tagstructbegin{tag=Title}%
    {\LARGE\bfseries\href{https://gma.math.ufl.edu/exams/topology-phd/}{Topology PhD Exam}, May 2017\par}%
  \tagstructend%
  \vskip 0.25em%
  \tagmcbegin{artifact}%
    \hrule height 1pt%
  \tagmcend%
  \vskip 1em%
}
\makeatother
\title{Topology PhD Exam, May 2017}
\author{}
\date{}
\begin{document}
\maketitle
\thispagestyle{empty}
\begin{examinstructions}
Work the following problems and show all work. Support all statements to the best of your ability.
\end{examinstructions}
\begin{examproblems}{0}
\item
Prove that a compact Hausdorff space is regular.
\item
Prove that if \(n>0\), every map \(S^n\to S^n\) is homotopic to a map with a fixed point.
\item
The oriented surface \(M_g\) of genus~\(g\), embedded in \(\mathbb{R}^3\) in the standard way, bounds a compact region~\(R\). Two copies of~\(R\), glued together by the identity map between their boundary surfaces \(M_g\), form a closed 3-manifold~\(X\). Compute the homology groups of~\(X\) using the Mayer–Vietoris sequence.
\item
Let \(f:S^1\to\mathbb{R}\) be a continuous map. Prove that there exists a point \(x\in S^1\) with \(f(x)=f(-x)\). (Note: do not simply state the Borsuk–Ulam Theorem; give a direct proof.)
\item
Let~\(M\) be a closed simply-connected orientable 3-manifold. Compute the integral homology and cohomology of~\(M\). What can you say about \(\pi_i(M)\), \(i\leq 3\)?
\end{examproblems}
\begin{examinstructions}
Answer the following with complete definitions, statements, or short proofs.
\end{examinstructions}
\begin{examproblems}{5}
\item
Prove that for a finite CW-complex~\(X\), \(H^1(X;\mathbb{Z})\) is torsion-free.
\item
Compute \(\chi(\mathbb{CP}^3\times\mathbb{RP}^2\times S^2)\).
\item
Give an example of a space that is path-connected but not locally path-connected.
\item
State the Urysohn Lemma.
\item
Prove that if \(m\neq n\), then \(\mathbb{R}^m\) is not homeomorphic to \(\mathbb{R}^n\).
\item
Does the function \(f:\mathbb{R}\to\mathbb{R}\) defined by \(f(x)=\arctan x\) admit a continuous extension \(\bar f:\beta\mathbb{R}\to\mathbb{R}\) to the Stone–Čech compactification? What about the function \(g(x)=e^x\)?
\item
State the Lefschetz Fixed Point Theorem.
\item
Describe all the connected covering spaces \(E\to S^1\).
\item
Does the following exact sequence of abelian groups necessarily split? Prove or give a counterexample. 
\[
0\to\mathbb{Z}\to A\to B\to 0
\]
\item
Compute the integral homology of the space \(\mathbb{RP}^4\times\mathbb{CP}^2\).
\end{examproblems}
\end{document}
